%!TEX root = ../main.tex
\chapter{Introduzione}
\label{cap:introduzione}

\section{Motivazioni}
\label{sec:intro-motivazioni}

Un incrocio semaforizzato stabilisce periodicamente quali flussi veicolari
autorizzare e quali arrestare. Su una singola intersezione il problema è
garantire che i veicoli in attesa continuino a scorrere in modo regolare.
Le soluzioni classiche affrontano questo compito in due modi. Un piano a
tempi fissi ripete una sequenza di fasi di durata prestabilita, calcolata
una tantum su stime di traffico medie e mai più rivista in tempo reale. Un
controllo attuato dai sensori concede invece il verde a una direzione solo
quando un sensore rileva la presenza di un veicolo in coda, adattandosi
così, entro certi limiti, alle variazioni del traffico osservato.

Se si considera una rete composta da più incroci, il problema cambia
natura. La fase scelta a un'intersezione influenza, con un ritardo
proporzionale alla distanza, il volume di veicoli in arrivo alle
intersezioni adiacenti, per cui le decisioni prese nei singoli punti della
rete non possono più essere indipendenti tra loro. Sistemi come SCOOT e
SCATS rispondono a questa esigenza da diversi decenni, aggiustando
periodicamente i piani semaforici di un'intera area urbana sulla base del
traffico rilevato in tempo reale \citep{wei2020tscsurvey}. Il controllo
semaforico su rete si configura quindi come un problema di coordinamento
distribuito, non di sola ottimizzazione isolata.

La diffusione di questi sistemi più adattivi nelle città resta comunque
parziale, più per ostacoli economici e infrastrutturali che per assenza
di alternative valide. Il problema del coordinamento semaforico resta
però un banco di prova interessante a prescindere dalla sua adozione
immediata sul campo, perché la stessa struttura, una rete di decisori
locali le cui azioni si propagano ai nodi vicini con un ritardo temporale,
ricorre in domini molto diversi dal traffico stradale. Lo stesso tipo di
coordinamento serve, per esempio, a instradare i pacchi nei punti di
smistamento di un centro logistico, a far muovere senza collisioni una
flotta di veicoli automatizzati in un magazzino, o a gestire un gruppo di
ascensori in un edificio in modo che nessun piano resti scoperto troppo a
lungo.

Valutare un controllore semaforico richiede di tenere conto di due
proprietà distinte. La prima è l'efficienza, la capacità di ridurre il
tempo che i veicoli impiegano ad attraversare la rete. La seconda è
l'equità, la garanzia che nessuna direzione di traffico venga penalizzata
sistematicamente a vantaggio delle altre. Le due proprietà non sono la
stessa cosa e non migliorano necessariamente insieme. Un controllore può
massimizzare il flusso complessivo e allo stesso tempo lasciare una
minoranza di veicoli bloccata a lungo, se la sua logica di decisione non
include alcuna attenzione esplicita al caso peggiore.

Tra i controllori reattivi allo stato dell'arte, \textsc{Max-Pressure}
\citep{varaiya2013maxpressure} è tra i più studiati. È una regola locale
che seleziona a ogni istante la fase capace di liberare più veicoli di
quanti ne farebbe accumulare, calcolata dalla sola differenza tra i
veicoli in coda sulle corsie in ingresso e in uscita che quella fase
servirebbe. Il risultato teorico che lo rende celebre riguarda proprio
l'orizzonte lungo, non solo la decisione istantanea. Se esiste una
qualunque strategia capace di impedire che le code di una rete crescano
senza limite nel tempo, \textsc{Max-Pressure} la trova, pur usando solo
informazione locale e senza conoscere in anticipo il volume di traffico
in arrivo. Si tratta però di un risultato dimostrato sotto ipotesi
teoriche idealizzate e verificato quasi soltanto in simulazione. Solo di
recente alcuni progetti pilota hanno iniziato a testarne il funzionamento
su hardware semaforico reale, e la sua diffusione resta comunque lontana
da quella di sistemi commerciali consolidati come SCOOT e SCATS. La sua
natura puramente greedy ha inoltre un costo sul fronte dell'equità.
Un'intersezione posta su un'arteria ad alto scorrimento tende a favorire
sistematicamente la direzione dominante, poiché la regola sceglie sempre
la fase con il valore istantaneo più alto, senza mai considerare da
quanto tempo attende la direzione esclusa. È una criticità che questo
lavoro osserva empiricamente nel Capitolo~\ref{cap:confronto}.

Alcuni metodi recenti in letteratura affrontano il problema penalizzando
esplicitamente le fasi responsabili di attese eccessive. FELight, per
esempio, introduce una metrica di equità basata sui tempi di attesa dei
veicoli e penalizza le fasi già identificate come responsabili di un
ritardo superiore a una soglia prestabilita \citep{xiao2025tscreview}.
Questa tesi persegue lo stesso obiettivo con un meccanismo strutturalmente
diverso, formalizzato nel Capitolo~\ref{cap:tsc-modelli}. Un termine
dedicato della ricompensa penalizza l'attesa massima osservata su una
corsia non servita dalla fase corrente, applicato a ogni singolo passo di
decisione e per ogni intersezione, non soltanto quando un ritardo supera
una soglia o riguarda una fase individuata in anticipo. È il motivo per
cui questo lavoro accetta, entro certi limiti, un tempo di percorrenza
medio peggiore pur di ottenere un'attesa massima più contenuta. A parità
di rete, un controllore che fa aspettare pochi veicoli molto a lungo non
è preferibile a uno leggermente più lento in media ma che non lascia mai
nessuno bloccato per un tempo indefinito.

Per verificare se questo meccanismo funzioni davvero, questo lavoro non
si affida al solo segnale di ricompensa, che dipende dai propri stessi
parametri e non è quindi un metro di giudizio imparziale, ma a una
metrica puramente osservazionale calcolata identicamente per ogni
controllore, indipendentemente da come è stato addestrato o se non è
stato addestrato affatto. La metrica è il tempo massimo che un veicolo
qualunque trascorre sulla stessa corsia senza superare l'incrocio, in una
qualunque intersezione della rete, definita formalmente nel
Capitolo~\ref{cap:confronto}.

L'obiettivo di questo lavoro è dunque duplice, e i suoi due termini non
sono sempre allineati. Da un lato ridurre il tempo di percorrenza medio e
massimo dei veicoli, l'\textbf{efficienza}. Dall'altro impedire che una
direzione di traffico venga sistematicamente sacrificata a vantaggio
delle altre, l'\textbf{equità}. Il Capitolo~\ref{cap:confronto} mostra
come questo compromesso si comporti concretamente al variare del peso
assegnato alla componente di equità nella ricompensa.

A questi due obiettivi se ne affianca un terzo, di natura diversa dai
primi due e altrettanto centrale in questo lavoro. Verificare che le
prestazioni misurate non dipendano dalla specifica rete stradale su cui
il modello è stato addestrato, ma si mantengano, almeno in parte, su
topologie, spaziature e condizioni di traffico mai incontrate durante il
training. Efficienza ed equità sono gli obiettivi della politica di
controllo. La \textbf{generalizzazione} è la condizione sotto cui questo
lavoro pretende che entrambi restino validi.

Gran parte della letteratura recente sul controllo semaforico appreso non
affronta sistematicamente questa domanda. La prassi più diffusa è
addestrare e valutare un modello sulla stessa rete stradale, al più
ripetendo l'intero esperimento, addestramento incluso, su una rete
diversa, non trasferendo un singolo modello già addestrato a una rete che
non ha mai visto \citep{wei2020tscsurvey}. CoLight \citep{wei2019colight}
conduce cinque esperimenti indipendenti, due su reti sintetiche e tre su
reti reali, Manhattan, Hangzhou e Jinan, ciascuno con il proprio
addestramento. MetaSTGAT \citep{wang2022metastgat}, il modello di
riferimento di questo lavoro, segue lo stesso schema su una griglia
sintetica e due dataset reali. Quando in questi lavori compare il termine
\emph{generalization}, indica quasi sempre la tenuta a diverse intensità
di traffico sulla stessa rete, non a una topologia strutturalmente
diversa \citep{wang2026astgcndrl, sun2026tgmaddpg}. Zhang et al. lo
rendono esplicito nella motivazione di GeneraLight
\citep{zhang2020generalight}. La maggior parte dei modelli di
reinforcement learning per il controllo semaforico è addestrata e
valutata nello stesso ambiente di traffico, il che porta a overfitting e
a una scarsa capacità di generalizzazione quando le condizioni reali
cambiano.

Le eccezioni che esistono restano circoscritte a un'unica dimensione di
variazione. AttendLight \citep{oroojlooy2020attendlight} addestra su un
insieme eterogeneo di intersezioni singole e ne verifica le prestazioni
su intersezioni mai viste, ma dichiara che l'estensione a una rete di più
intersezioni coordinate resta lavoro futuro. Advanced-XLight
\citep{zhang2022advancedxlight} misura un rapporto di trasferimento tra
un modello addestrato su un dataset reale e valutato direttamente su un
altro, senza ulteriore addestramento, ma sempre tra reti di scala
paragonabile. Lo stesso GeneraLight generalizza al variare del flusso di
traffico, non della topologia della rete.

Questo lavoro adotta invece la generalizzazione come requisito
sistematico su assi che la letteratura esaminata non copre
congiuntamente. Ogni modello, sia di reinforcement learning sia di
riferimento classico, è addestrato una sola volta su una griglia
$4\times4$ e valutato, senza alcun riaddestramento, su topologie di
dimensione maggiore mai incontrate, $5\times5$ e $6\times6$, su una
spaziatura stradale diversa, e su profili di traffico differenti da
quello di addestramento. Il Capitolo~\ref{cap:confronto} descrive nel
dettaglio il protocollo sperimentale adottato.

Il controllo appreso tramite reinforcement learning propone un paradigma
diverso da quello di \textsc{Max-Pressure}. Invece di una regola scritta
a mano, una politica $\pi$ viene addestrata a massimizzare una ricompensa
che codifica l'obiettivo di rete, osservando ripetutamente le conseguenze
delle proprie azioni in simulazione. Applicato a una singola
intersezione, il problema è modellabile come un processo decisionale di
Markov standard. Su scala di rete assume invece la forma di un sistema
multi-agente a osservabilità parziale, dove ogni nodo rileva
esclusivamente il proprio intorno locale, ed emerge quindi la necessità
di un meccanismo per la condivisione delle informazioni tra agenti
adiacenti. Le reti neurali a grafo assolvono a questa funzione.
Rappresentando la rete stradale come un grafo, con le intersezioni come
nodi e le strade come archi, l'informazione si propaga tra nodi vicini,
permettendo a ogni intersezione di condizionare la propria decisione
sullo stato dei suoi vicini senza richiedere una politica centralizzata.

\textbf{MetaSTGAT} \citep{wang2022metastgat} combina questi aspetti con
uno ulteriore. Invece di condividere gli stessi pesi tra tutte le
intersezioni della rete, un modulo di \emph{meta-learning} genera
dinamicamente i pesi di uno strato di attenzione e di uno strato
ricorrente a partire da caratteristiche locali dell'intersezione, così
che nodi diversi possano essere serviti da trasformazioni diverse pur
restando un unico modello addestrato end-to-end. Gli autori originali
riportano una riduzione del travel time medio rispetto a CoLight
\citep{wei2019colight}, un metodo di riferimento basato su \emph{graph
attention} ma con pesi condivisi. La combinazione di \emph{attention} su
grafo, ricorrenza e reinforcement learning non è isolata. Varianti di
questo framework dominano la letteratura recente
\citep{wang2026astgcndrl, sun2026tgmaddpg, hu2026mpnnlight}.

Fin qui si è discusso come il controllo semaforico su rete sia un
problema di coordinamento distribuito, in cui l'efficienza aggregata non
implica alcuna garanzia di equità tra le direzioni di traffico servite,
come dimostra il caso di \textsc{Max-Pressure}. I meccanismi di equità
già proposti in letteratura intervengono in modo selettivo, solo oltre
soglie o su fasi identificate in anticipo, mentre le architetture di
controllo appreso capaci di coordinare intere reti di intersezioni, come
MetaSTGAT, offrono il meccanismo di coordinamento ma non trattano
l'equità come obiettivo esplicito di progettazione. La generalizzazione a
topologie di rete, spaziature stradali o profili di traffico mai
osservati in addestramento resta infine raramente verificata in modo
sistematico, anche quando il termine compare esplicitamente nei lavori
citati. Questa tesi affronta congiuntamente questi tre aspetti,
replicando da zero l'architettura e la procedura di addestramento di
MetaSTGAT come base sperimentale.

\section{Obiettivi e Contributi di questo lavoro}

L'obiettivo principale di questo lavoro è verificare se un controllore
appreso e coordinato tramite reti neurali a grafo possa garantire equità
tra direzioni di traffico senza rinunciare all'efficienza raggiunta da un
controllore puramente reattivo come \textsc{Max-Pressure}, e se tale
garanzia si mantenga su topologie, spaziature e profili di traffico mai
incontrati in addestramento. Un obiettivo secondario, strumentale al
primo, è isolare quale parte dell'architettura e della procedura di
addestramento sia effettivamente responsabile del comportamento
osservato, invece di attribuirlo indistintamente all'insieme del modello.
Nello specifico, i contributi di questa tesi sono i seguenti.

\paragraph{Un ambiente di simulazione multi-intersezione riproducibile.}
Questa tesi costruisce da zero un ambiente su \textsc{CityFlow}
\citep{zhang2019cityflow}, con reti stradali a griglia di dimensione
$4\times4$, $5\times5$ e $6\times6$, due spaziature stradali distinte,
otto fasi semaforiche per intersezione e una formalizzazione del
problema dettagliata nel Capitolo~\ref{cap:tsc-modelli}, insieme a una
generazione parametrica dei dati di traffico con profili a intensità
costante, a gradino e a giornata lavorativa. In ogni configurazione il
traffico su una riga della griglia è potenziato artificialmente per
simulare un'arteria stradale, condizione necessaria per poter osservare
uno squilibrio direzionale quando il traffico non è simmetrico.

\paragraph{Uno studio sistematico del compromesso tra efficienza ed
equità.} Variando il peso assegnato alla componente di equità nella
ricompensa, questa tesi quantifica direttamente il compromesso che il
meccanismo introduce, mostrando in particolare che rimuoverlo del tutto
non produce un semplice guadagno di efficienza, ma un collasso della
garanzia di equità verso il comportamento di \textsc{Max-Pressure}, come
mostra il Capitolo~\ref{cap:confronto}.

\paragraph{Uno studio di ablation sull'origine del comportamento
osservato.} Invece di presentare le differenze rispetto alla procedura
originale come un insieme indistinto di miglioramenti, questa tesi isola
singolarmente ognuna di esse, dal meccanismo di stato e ricompensa fino ai
dettagli dell'algoritmo di apprendimento, per stabilire quale parte sia
effettivamente necessaria alla garanzia di equità osservata e quale
sia invece indipendente da essa.

\paragraph{Un confronto tra meccanismi di comunicazione spaziale a parità
di resto dell'architettura.} Questa tesi confronta tre meccanismi di
comunicazione tra intersezioni, \emph{graph attention}, convoluzione
grafica e propagazione a onda ispirata a SONAR, mantenendo identico il
resto dell'architettura, per verificare se il comportamento osservato
dipenda specificamente dal meccanismo di attenzione appreso o sia una
proprietà più generale della coordinazione locale via rete neurale a
grafo.

\paragraph{Un protocollo di valutazione a generalizzazione sistematica.}
Ogni modello di questa tesi, appreso o classico, è valutato senza
riaddestramento su topologie, spaziature e profili di traffico mai
incontrati in addestramento, riportando media e caso peggiore del tempo
di viaggio, la percentuale di veicoli che completano il proprio percorso
e la massima attesa registrata da un veicolo allo stesso incrocio.

\section{Struttura della tesi}

Il Capitolo~\ref{cap:background} introduce il background necessario a
leggere i capitoli successivi senza assumere familiarità pregressa con
nessuno dei campi coinvolti, dal reinforcement learning alle reti neurali
ricorrenti, dalle reti neurali a grafo al controllo semaforico. Il
Capitolo~\ref{cap:tsc-modelli} formalizza il problema affrontato,
descrive l'ambiente di simulazione e le baseline classiche, e presenta
l'architettura di MetaSTGAT e le sue varianti insieme a ogni differenza
introdotta rispetto all'articolo originale, motivata caso per caso. Il
Capitolo~\ref{cap:confronto} descrive il protocollo sperimentale e
riporta i risultati ottenuti sui tre studi, peso dell'equità, ablation,
meccanismi spaziali. Il Capitolo~\ref{cap:conclusioni} discute i
risultati alla luce dei tre obiettivi dichiarati in questa introduzione,
indica le direzioni lasciate aperte e dichiara i limiti dello studio.
